%% ===============================
%% anlagen/listings.tex
%% Quellcode-Auszüge: Hackintosh Companion App
%% ===============================

% ===============================
% Listings Setup für Dart/Flutter
% ===============================
\lstdefinelanguage{Dart}{
    morekeywords={abstract,as,assert,async,await,break,case,catch,class,const,continue,covariant,default,deferred,do,dynamic,else,enum,export,extends,extension,external,factory,false,final,finally,for,Function,get,hide,if,implements,import,in,interface,is,late,library,mixin,new,null,on,operator,part,rethrow,return,set,show,static,super,switch,sync,this,throw,true,try,typedef,var,void,while,with,yield},
    sensitive=true,
    morecomment=[l]{//},
    morecomment=[s]{/*}{*/},
    morestring=[b]',
    morestring=[b]"
}

\lstset{
    language=Dart,
    basicstyle=\ttfamily\small,
    numbers=left,
    numberstyle=\tiny\color{gray},
    stepnumber=1,
    numbersep=5pt,
    frame=single,
    rulecolor=\color{black!30},
    backgroundcolor=\color{white!97!gray},
    breaklines=true,
    captionpos=b,
    commentstyle=\color{green!50!black},
    keywordstyle=\color{blue},
    stringstyle=\color{orange}
}

% ===============================
% Datenbank-Logik (SQLite)
% ===============================
\section{Datenbank-Implementierung (SQLite)}
Im folgenden Listing wird die Initialisierung der lokalen Datenbank sowie das relationale Schema mit Fremdschlüssel-Abhängigkeiten dargestellt.

\begin{lstlisting}[language=Dart, caption={DatabaseHelper.dart - Relationales Schema und Initialisierung}]
// Auszug aus lib/database/database_helper.dart
Future _createDB(Database db, int version) async {
  // Haupttabelle fuer Geraete
  await db.execute('''
    CREATE TABLE devices (
      id INTEGER PRIMARY KEY AUTOINCREMENT,
      name TEXT NOT NULL,
      cpu TEXT NOT NULL,
      gpu TEXT NOT NULL,
      compatible INTEGER NOT NULL
    )
  ''');

  // Tabelle fuer Kexts mit Relation zu devices
  await db.execute('''
    CREATE TABLE kexts (
      id INTEGER PRIMARY KEY AUTOINCREMENT,
      device_id INTEGER,
      name TEXT NOT NULL,
      version TEXT,
      FOREIGN KEY (device_id) REFERENCES devices (id) ON DELETE CASCADE
    )
  ''');
}
\end{lstlisting}

% ===============================
% REST-API Synchronisation
% ===============================
\section{REST-API Synchronisation}
Dieses Listing zeigt die Implementierung des Offline-First-Ansatzes durch den Abgleich lokaler SQLite-Daten mit einem remote JSON-Endpunkt.

\begin{lstlisting}[language=Dart, caption={SyncService.dart - REST-API Anbindung}]
// Auszug aus lib/services/sync_service.dart
Future<void> syncUp() async {
  final localDevices = await DatabaseHelper.instance.getAllDevices();
  
  for (var device in localDevices) {
    try {
      await http.post(
        Uri.parse('$baseUrl/devices'),
        body: json.encode(device.toMap()),
        headers: {'Content-Type': 'application/json'},
      );
    } catch (e) {
      debugPrint("Upload-Fehler: $e");
    }
  }
}
\end{lstlisting}

% ===============================
% 3D-Integration
% ===============================
\section{3D-Visualisierung}
Auszug der Widget-Logik zur Einbindung des interaktiven Laptop-Modells im Detail-Screen.

\begin{lstlisting}[language=Dart, caption={VisualScreen.dart - 3D-Modell Einbindung}]
// Auszug aus lib/screens/visual_screen.dart
@override
Widget build(BuildContext context) {
  return Flutter3DViewer(
    controller: controller,
    src: 'assets/models/thinkpad.glb', // Pfad zum 3D-Asset
  );
}
\end{lstlisting}

% ===============================
% Boot-Animation
% ===============================
\section{Simulation der Boot-Sequenz}
Implementierung der zeitgesteuerten Phasen (BIOS, OpenCore, macOS) zur Steigerung der User Experience.

\begin{lstlisting}[language=Dart, caption={BootAnimationScreen.dart - Phasensteuerung}]
// Auszug aus lib/screens/boot_animation_screen.dart
void _startBootSequence() async {
  // Phase 1: BIOS Logs
  await Future.delayed(Duration(milliseconds: 300));
  setState(() => _bootPhase = 0);

  // Phase 2: OpenCore Picker
  await Future.delayed(Duration(seconds: 2));
  setState(() => _bootPhase = 1);

  // Phase 3: macOS Start
  await Future.delayed(Duration(seconds: 2));
  setState(() => _bootPhase = 2);
}
\end{lstlisting}